% =====================================================================
%  FINAL — Methodology section for the Semantic Analysis chapter (BERTopic)
%  Figures verified against:
%    steam-data/step03-output/{en,pt}/optuna/best_params.json
%    steam-data/step03-output/{en,pt}/stability/stability_report.json
%    step03_bertopic/{METHODOLOGY_REPORT.md, RESULTS_REPORT.md}
%  Edits vs. previous draft (see chat): fixed $c_{npmi}$ math typo,
%  fixed "hyperparameters, of ... jointly" punctuation, and added the
%  English final-model fit statistics for parallelism with Portuguese.
% =====================================================================

\section{Methodology}
\label{sec:semantic-method}

Topic modeling pipeline was built with BERTopic\footnote{BERTopic official documentation. \url{https://maartengr.github.io/BERTopic/}. Accessed: July 20, 2026.} \cite{grootendorst:2022}, to discover latent thematic clusters within the toxic corpus. To ensure structural consistency while respecting linguistic nuances of each dataset, the English and Portuguese corpora were processed with the same underlying code, independently, through an identical seven-stage methodology.

To guarantee reproducibility, a global seed ($seed = 42$) was fixed across every stochastic stage of the pipeline, namely UMAP's fitting, HDBSCAN's clustering, the Optuna search itself, the incremental subsampling performed during the stability analysis described below, and the random sampling of reviews later shown to the labeling model. Parallelism was disabled in both UMAP and HDBSCAN, since concurrent execution would otherwise introduce non-determinism beyond the reach of a fixed seed alone.

Reviews were first cleaned following the same rigor of preprocessing described in other parts of this study---e.g., Chapter~\ref{chap:charact-tfidf}---lowercasing text, stripping URLs, normalizing accented characters through NFKD normalization to ASCII, and discarding non-alphabetic characters. Cleaned reviews were then converted into dense vector representations through the pretrained model \texttt{paraphrase-multilingual-MiniLM-L12-v2}, distributed via the Sentence-Transformers library\footnote{Sentence-Transformers official documentation. \url{https://www.sbert.net/}. Accessed: July 20, 2026.}. This model was chosen for producing a shared multilingual embedding space, positioning semantically similar reviews close together in vector space regardless of whether they were written in English or Portuguese, which allows the same architecture to be applied uniformly to both language subsets. Because direct clustering over these high-dimensional embeddings is computationally prohibitive, and density-based methods in particular tend to degrade in high-dimensional spaces, a consequence of the curse of dimensionality \cite{bellman:1957,beyer:1999}, embeddings were immediately compressed with PCA to 50 dimensions, preserving the most significant share of global variance as a preliminary reduction.

The core clustering algorithm required precise calibration to yield meaningful semantic groupings. The Optuna\footnote{Optuna official documentation. \url{https://optuna.org/}. Accessed: July 20, 2026.} framework \cite{akiba:2019} was employed to jointly optimize the hyperparameters of both UMAP and HDBSCAN, across 60 trials per language. The optimization target combined two competing signals into a single objective to be minimized, the fraction of documents left as outliers (topic $-1$), minus the $c_{npmi}$ topic coherence measured over each topic's top 10 words, so that a configuration was rewarded for simultaneously lowering the outlier fraction and raising coherence, rather than for maximizing coherence alone. To prevent the algorithm from collapsing the corpus into an uninformative handful of massive clusters, trials producing fewer than 10 topics were heavily penalized, a floor set low enough not to bias the search toward an artificially fine-grained partition, while still ruling out the degenerate two- or three-topic configurations the corpus's scale would otherwise allow.

The original search range for HDBSCAN's minimum cluster size, $[100, 3000]$, was set as an initial heuristic estimate, and carried no explicit relationship to corpus size. This became a problem once the pipeline was applied to both languages, since the same absolute range, adequate for the $480{,}729$-review English corpus, represented a proportionally much larger share of the smaller Portuguese corpus, structurally biasing its search toward degenerate, few-topic solutions. The search space was therefore re-expressed as a fraction of each language's corpus size, ranging from $0.0208\%$ to $0.6241\%$, calibrated to reproduce, exactly, the original absolute range when applied to the English corpus ($100/480{,}729 \approx 0.0208\%$; $3{,}000/480{,}729 \approx 0.6241\%$), so that English search behavior was left unchanged while Portuguese search behavior became proportionally equivalent to it.

Complementarily, an empirical stability analysis was conducted to determine how large a training sample needs to be for the resulting topic structure to stop changing meaningfully as more data is added. This was assessed by fitting the c-TF-IDF representation over the complete corpus once, so that every sample size compared shares the same vocabulary and column space, and comparing the cosine similarity between the topic-term matrices obtained at successive sample sizes, evaluated across incremental subsets of the corpus, namely $10\%$, $20\%$, $30\%$, and $40\%$ fractions. A sample size was accepted as sufficient once this similarity crossed a threshold of $0.85$ against the preceding, smaller size, and its resulting topic count also cleared the same ten-topic floor used during hyperparameter search, since two consecutive sample sizes can otherwise collapse onto the same degenerate, near-empty structure and register as artificially stable simply by being equally uninformative.

This stability analysis, however, only establishes that the topic structure itself converges as sample size grows, not that a model trained on a partial sample generalizes, through inference, to the documents left out of it. The two languages diverged on this second question. For Portuguese, both held: structure converged at $40\%$ of the corpus, with a cosine similarity of $0.9079$ against the $30\%$ step, and classifying the full corpus with a model trained on that $40\%$ sample populated every recommended topic correctly. The final Portuguese model was accordingly trained on $13{,}474$ reviews, yielding 15 topics, with a fit-time outlier rate of $29.50\%$ and a $c_{npmi}$ coherence of $-0.0638$.

For English, structure likewise converged at a partial sample, specifically the $30\%$ step, but classifying the remaining, much larger, held-out portion of the corpus with a model trained on it did not generalize. When a model trained on the $144{,}219$-review search sample was used to classify the remaining, unseen English reviews, essentially all of that held-out volume collapsed into only two of the twenty topics found during that training run, leaving the other eighteen without a single new document. This could indicate that UMAP's extrapolation to unseen points becomes unreliable once the volume being extrapolated exceeds the volume it was fitted on by roughly an order of magnitude, an order of magnitude reached in English but not in Portuguese. The final English model was therefore trained on the complete toxic corpus of $480{,}729$ reviews instead, removing the need for any such extrapolation at inference time, since every classified document was already seen during training, and yielded 23 topics, with a fit-time outlier rate of $25.00\%$ and a $c_{npmi}$ coherence of $0.0498$. For both languages, each topic is additionally represented by the ten terms that best distinguish it from every other topic identified in the corpus, computed via c-TF-IDF.

With the models trained, the pipeline proceeded to the inference and consolidation stage. The entire toxic corpus for each language, $480{,}729$ English reviews and $33{,}685$ Portuguese reviews, was classified to assign a topic label to every review, leaving an outlier rate of $4.03\%$ in English and $4.23\%$ in Portuguese.

The final stage involved a semantic labeling assisted by an LLM, \texttt{claude-fable-5} through the Anthropic API\footnote{Anthropic API documentation. \url{https://docs.claude.com/}. Accessed: July 20, 2026.}. For each non-outlier topic, a sample of 20 reviews was extracted directly from the classified corpus, under a fixed seed, alongside the topic's top keywords from c-TF-IDF. The model was prompted to assign a concise title, a brief description, and a category from a predefined, closed taxonomy, namely \texttt{generic\_insult}, \texttt{violence\_threat}, \texttt{sexual\_content}, \texttt{identity\_hate}, \texttt{humor\_sarcasm}, \linebreak\texttt{nationality\_xenophobia}, \texttt{product\_complaint}, and \texttt{other}. The full prompt used for this stage is provided in Appendix~\ref{app:semantic-prompt}. Outlier clusters (topic $-1$) were deliberately excluded from this LLM inference and cataloged separately as unclustered noise.
